% ============================================================
%  CHAPTER 6 — IMPLEMENTATION
% ============================================================
\chapter{Project Implementation}


\section{Deep Learning Model Training}

\subsection{Dataset Preparation}

The \textbf{RAF-DB (Real-world Affective Faces Database)} dataset was used to
train the emotion recognition model. It contains 15,339 facial images
annotated with seven basic emotion categories: Angry, Disgust, Fear, Happy,
Neutral, Sad, and Surprise. The dataset was split into 12,271 training images
and 3,068 test images following the official partition.

Preprocessing steps applied to each image:
\begin{itemize}
    \item Resize to $224 \times 224$ pixels (ResNet-18 input size).
    \item Random horizontal flip (probability 0.5) for data augmentation.
    \item Random rotation in the range $[-10°, +10°]$.
    \item Colour jitter (brightness, contrast, saturation variation of 0.2).
    \item Normalisation using ImageNet mean $[0.485, 0.456, 0.406]$ and
          standard deviation $[0.229, 0.224, 0.225]$.
\end{itemize}

Class imbalance in the dataset (Disgust and Fear are underrepresented) was
addressed using \textbf{weighted random sampling} during training, where the
sampling weight of each class is inversely proportional to its frequency.

\subsection{Model Architecture}

\textbf{ResNet-18} (Residual Network with 18 layers) was selected as the
base architecture. It was pretrained on ImageNet and fine-tuned for the
7-class emotion classification task. The key modification was replacing the
final fully-connected layer (originally 1000 output units for ImageNet) with
a new linear layer of 7 output units:

\begin{itemize}
    \item \textbf{Input}: $3 \times 224 \times 224$ RGB image tensor
    \item \textbf{Feature Extractor}: ResNet-18 convolutional backbone
          (frozen for first 3 epochs, then fully fine-tuned)
    \item \textbf{Classifier Head}: Linear(512, 7) + Softmax
\end{itemize}

\subsection{Training Configuration}

\begin{table}[H]
\centering
\caption{Model Training Hyperparameters}
\begin{tabular}{ll}
\toprule
\textbf{Hyperparameter} & \textbf{Value} \\
\midrule
Optimiser               & Adam \\
Initial learning rate   & $1 \times 10^{-3}$ \\
LR scheduler            & StepLR (step size = 10, gamma = 0.1) \\
Batch size              & 64 \\
Epochs                  & 30 \\
Loss function           & CrossEntropyLoss \\
Weight decay            & $1 \times 10^{-4}$ \\
Early stopping patience & 5 epochs \\
\bottomrule
\end{tabular}
\end{table}

Training was performed on a GPU-accelerated environment (NVIDIA CUDA).
The best model checkpoint (highest validation accuracy) was saved and
exported as a \texttt{.pth} file for deployment.

\subsection{Model Performance}

The trained ResNet-18 model achieved the following results on the RAF-DB test set:

\begin{table}[H]
\centering
\caption{Model Evaluation Results on RAF-DB Test Set}
\begin{tabular}{lcc}
\toprule
\textbf{Emotion}  & \textbf{Precision} & \textbf{Recall} \\
\midrule
Angry             & 0.81               & 0.78 \\
Disgust           & 0.74               & 0.70 \\
Fear              & 0.76               & 0.73 \\
Happy             & 0.93               & 0.95 \\
Neutral           & 0.88               & 0.89 \\
Sad               & 0.82               & 0.80 \\
Surprise          & 0.87               & 0.86 \\
\midrule
\textbf{Overall Accuracy} & \multicolumn{2}{c}{\textbf{85.7\%}} \\
\bottomrule
\end{tabular}
\end{table}

\section{Backend Implementation}

\subsection{Project Structure}

The FastAPI backend is organised as follows:

\begin{itemize}
    \item \texttt{main.py} — FastAPI application entry point; registers
          routers and CORS middleware.
    \item \texttt{routers/auth.py} — Authentication endpoints
          (\texttt{/register}, \texttt{/login}).
    \item \texttt{routers/analysis.py} — Inference endpoint
          (\texttt{/api/analyze}) and session management.
    \item \texttt{services/emotion\_service.py} — MediaPipe and
          ResNet-18 inference logic.
    \item \texttt{services/engagement\_service.py} — Head pose and
          blink-rate computation.
    \item \texttt{services/db\_service.py} — Async MongoDB operations.
    \item \texttt{models/schemas.py} — Pydantic request/response models.
    \item \texttt{core/security.py} — JWT and bcrypt utilities.
\end{itemize}

\subsection{Emotion Inference Endpoint}

The core inference endpoint accepts a Base64-encoded image frame and returns
the emotion prediction and engagement metrics:

\begin{enumerate}
    \item The Base64 payload is decoded to a NumPy array using OpenCV.
    \item MediaPipe Face Mesh processes the frame to extract 468 facial
          landmarks and compute a face bounding box.
    \item The face region is cropped, resized to $224 \times 224$, and
          normalised before being passed to the ResNet-18 model.
    \item The model returns a softmax probability vector; the argmax class
          is selected as the predicted emotion.
    \item In parallel, head pose is estimated from 6 key landmarks using
          \texttt{cv2.solvePnP}, and blink rate is computed from the
          Eye Aspect Ratio (EAR) of both eyes.
    \item The composite engagement score is computed and the full result is
          serialised as a JSON response.
\end{enumerate}

\subsection{Authentication Implementation}

User registration stores the bcrypt-hashed password (work factor 12) in
MongoDB. Login verifies the password hash and issues a signed JWT token with
a configurable expiry (default: 60 minutes). All protected API routes use
a FastAPI dependency (\texttt{get\_current\_user}) that decodes and validates
the token before passing the user context to the handler.

\subsection{CORS and Deployment Configuration}

Cross-Origin Resource Sharing (CORS) middleware is configured to allow
requests from the Vercel-hosted frontend domain. The application is
containerised using Docker with a \texttt{uvicorn} ASGI server and deployed
to Hugging Face Spaces via a \texttt{Dockerfile} that installs all Python
dependencies and copies the trained model weights.

\section{Frontend Implementation}

\subsection{Project Structure}

The React + Vite frontend is structured using a feature-based layout:

\begin{itemize}
    \item \texttt{src/pages/} — Top-level page components (Login, Register,
          Dashboard, History, Admin).
    \item \texttt{src/components/} — Reusable UI components (WebcamFeed,
          EmotionOverlay, EngagementChart, SessionTable).
    \item \texttt{src/services/api.js} — Axios instance with JWT
          interceptor for attaching the Bearer token to every request.
    \item \texttt{src/store/} — Global state management using React Context
          for authentication and session state.
    \item \texttt{src/hooks/} — Custom hooks (\texttt{useWebcam},
          \texttt{useAnalysis}) for webcam access and API polling.
\end{itemize}

\subsection{Webcam Capture Module}

The \texttt{useWebcam} hook uses the \texttt{navigator.mediaDevices.getUserMedia}
Web API to access the camera stream. Frames are captured at a configurable
interval (200 ms default) by drawing the video stream onto a hidden
\texttt{<canvas>} element and calling \texttt{canvas.toDataURL('image/jpeg', 0.7)}
to produce a compressed Base64 string for transmission.

\subsection{Live Dashboard}

The main dashboard renders:
\begin{itemize}
    \item A live webcam feed with the detected emotion label and confidence
          overlaid as a transparent badge.
    \item An engagement score gauge updated in real time.
    \item A Recharts \texttt{LineChart} plotting engagement score over the
          session timeline.
    \item A \texttt{PieChart} showing the distribution of detected emotions
          for the current session.
\end{itemize}

Session start/stop controls trigger API calls to create and finalise session
records in MongoDB.

\section{Database Integration}

The backend uses \textbf{Motor}, the asynchronous Python driver for MongoDB,
to perform non-blocking database operations. A connection pool is initialised
at application startup and shared across all request handlers via FastAPI's
dependency injection system. Key operations include:

\begin{itemize}
    \item \textbf{User creation}: Insert a new document into the
          \texttt{users} collection on registration.
    \item \textbf{Session creation}: Insert a new session document on
          session start.
    \item \textbf{Frame append}: Use MongoDB's \texttt{\$push} operator to
          append each frame analysis record to the session's \texttt{frames}
          array atomically.
    \item \textbf{Session finalisation}: Update \texttt{end\_time},
          \texttt{avg\_engagement}, and \texttt{dominant\_emotion} on
          session end.
    \item \textbf{History retrieval}: Query sessions by \texttt{user\_id}
          with pagination using \texttt{skip} and \texttt{limit}.
\end{itemize}

\section{Deployment}

\subsection{Backend Deployment — Hugging Face Spaces}

The FastAPI backend is deployed as a Docker Space on Hugging Face. The
\texttt{Dockerfile} performs the following steps:

\begin{enumerate}
    \item Uses \texttt{python:3.10-slim} as the base image.
    \item Installs system dependencies (libGL for OpenCV).
    \item Copies and installs Python packages from \texttt{requirements.txt}.
    \item Copies the application source and pre-trained model weights.
    \item Exposes port 7860 and starts the Uvicorn server.
\end{enumerate}

Environment variables (MongoDB URI, JWT secret key) are configured through
the Hugging Face Spaces secrets interface to avoid hardcoding sensitive
credentials.

\subsection{Frontend Deployment — Vercel}

The React application is deployed on Vercel by connecting the GitHub
repository. Vercel automatically builds the Vite project on each push to the
main branch using the command \texttt{npm run build}. The backend API URL is
configured as a Vercel environment variable (\texttt{VITE\_API\_URL}) and
injected at build time.
